% 分野別類題: 微分計算（多変数の連鎖律・偏微分・全微分）
% z=f(x,y) の合成関数の連鎖律、極座標変換、陰関数定理を扱う。計算スピード重点分野のため5問（基本3問＋やや複雑2問）。
% コンパイル: TEXINPUTS="../../../00_common:$TEXINPUTS" latexmk -xelatex problems.tex
\documentclass[a4paper,11pt]{article}
\usepackage{preamble}

\begin{document}

\kamokumei{類題演習 --- 微分計算（多変数の連鎖律・偏微分・全微分）}

\shijibun{以下の問いに答えよ。必要ならば次の連鎖律（合成関数の微分公式）を用いてよい。
$z = f(x,y)$、$x = x(t)$、$y = y(t)$ のとき
\[
\frac{dz}{dt} = \frac{\partial z}{\partial x}\frac{dx}{dt} + \frac{\partial z}{\partial y}\frac{dy}{dt}.
\]
また $z=f(x,y)$、$x=x(u,v)$、$y=y(u,v)$ のときは
\[
\frac{\partial z}{\partial u} = \frac{\partial z}{\partial x}\frac{\partial x}{\partial u} + \frac{\partial z}{\partial y}\frac{\partial y}{\partial u}, \qquad
\frac{\partial z}{\partial v} = \frac{\partial z}{\partial x}\frac{\partial x}{\partial v} + \frac{\partial z}{\partial y}\frac{\partial y}{\partial v}.
\]
陰関数 $F(x,y)=0$ が定める $y=y(x)$ については $\dfrac{dy}{dx} = -\dfrac{F_x}{F_y}$ が成り立つ。}

\shomon{問1.}{$z = x^{2}y - y^{3}$、$x = e^{t}$、$y = \cos t$ のとき、合成関数の連鎖律を用いて $\dfrac{dz}{dt}$ を求めよ。}

\shomon{問2.}{$z = x^{2} - y^{2}$ について、極座標変換 $x = r\cos\theta$、$y = r\sin\theta$ を行うとき、$\dfrac{\partial z}{\partial r}$ と $\dfrac{\partial z}{\partial \theta}$ を求めよ。}

\shomon{問3.}{陰関数 $x^{3} + y^{3} - 3xy = 0$ について、一般の点における $\dfrac{dy}{dx}$ を $x,y$ の式で表せ。さらに、この曲線上の点 $\left(\dfrac{3}{2}, \dfrac{3}{2}\right)$ における $\dfrac{dy}{dx}$ の値を求めよ。}

\shomon{問4.}{$z = f(x,y)$ に対し、$x = u^{2} - v^{2}$、$y = 2uv$ とおくとき、$\dfrac{\partial z}{\partial u}$、$\dfrac{\partial z}{\partial v}$ を $\dfrac{\partial z}{\partial x}$、$\dfrac{\partial z}{\partial y}$、$u$、$v$ を用いて表せ。さらに $f(x,y) = xy$ の場合について、$\dfrac{\partial z}{\partial u}$、$\dfrac{\partial z}{\partial v}$ を $u,v$ の式として具体的に求めよ。}

\shomon{問5.}{$z = f(x,y)$ に対し、極座標変換 $x = r\cos\theta$、$y = r\sin\theta$ を行うとき、2階偏微分について次の関係式が成り立つことを示せ。
\[
\frac{\partial^{2} z}{\partial x^{2}} + \frac{\partial^{2} z}{\partial y^{2}} = \frac{\partial^{2} z}{\partial r^{2}} + \frac{1}{r}\frac{\partial z}{\partial r} + \frac{1}{r^{2}}\frac{\partial^{2} z}{\partial \theta^{2}}.
\]}

\end{document}
